
\documentclass{article}
\usepackage{colortbl}
\usepackage{makecell}
\usepackage{multirow}
\usepackage{supertabular}

\begin{document}

\newcounter{utterance}

\centering \large Interaction Transcript for game `hot\_air\_balloon', experiment `air\_balloon\_survival\_en\_complexity\_hard', episode 0 with Llama{-}3.3{-}70B{-}Instruct{-}t0.0.
\vspace{24pt}

{ \footnotesize  \setcounter{utterance}{1}
\setlength{\tabcolsep}{0pt}
\begin{supertabular}{c@{$\;$}|p{.15\linewidth}@{}p{.15\linewidth}p{.15\linewidth}p{.15\linewidth}p{.15\linewidth}p{.15\linewidth}}
    \# & $\;$A & \multicolumn{4}{c}{Game Master} & $\;\:$B\\
    \hline

    \theutterance \stepcounter{utterance}  
    & & \multicolumn{4}{p{0.6\linewidth}}{
        \cellcolor[rgb]{0.9,0.9,0.9}{
            \makecell[{{p{\linewidth}}}]{
                \texttt{\tiny{[P1$\langle$GM]}}
                \texttt{You are participating in a collaborative negotiation game.} \\
\\ 
\texttt{Together with another participant, you must agree on a single set of items that will be kept. Each of you has your own view of how much each item matters to you (importance). You do not know how the other participant values the items. Additionally, you are given the effort each item demands.} \\
\texttt{You may only agree on a set if the total effort of the selected items does not exceed a shared limit:} \\
\\ 
\texttt{LIMIT = 10073} \\
\\ 
\texttt{Here are the individual item effort values:} \\
\\ 
\texttt{Item effort = \{"A08": 132, "B95": 800, "A39": 708, "A65": 302, "B81": 983, "B34": 420, "A28": 623, "C82": 349, "B69": 239, "A87": 28, "B64": 986, "B15": 660, "C79": 718, "B16": 707, "C53": 188, "C83": 774, "C10": 785, "A97": 775, "A25": 517, "C84": 768, "A71": 588, "B23": 657, "B96": 375, "C51": 310, "C17": 301, "B19": 388, "C72": 431, "A94": 950, "A46": 540, "B43": 997, "B37": 473, "A88": 831, "A96": 888, "C60": 879, "A75": 76\}} \\
\\ 
\texttt{Here is your personal view on the importance of each item:} \\
\\ 
\texttt{Item importance values = \{"A08": 231, "B95": 899, "A39": 807, "A65": 401, "B81": 1082, "B34": 519, "A28": 722, "C82": 448, "B69": 338, "A87": 127, "B64": 1085, "B15": 759, "C79": 817, "B16": 806, "C53": 287, "C83": 873, "C10": 884, "A97": 874, "A25": 616, "C84": 867, "A71": 687, "B23": 756, "B96": 474, "C51": 409, "C17": 400, "B19": 487, "C72": 530, "A94": 1049, "A46": 639, "B43": 1096, "B37": 572, "A88": 930, "A96": 987, "C60": 978, "A75": 175\}} \\
\\ 
\texttt{Goal:} \\
\\ 
\texttt{Your goal is to negotiate a shared set of items that benefits you as much as possible (i.e., maximizes total importance to YOU), while staying within the LIMIT. You are not required to make a PROPOSAL in every message {-} you can simply negotiate as well. All tactics are allowed!} \\
\\ 
\texttt{Interaction Protocol:} \\
\\ 
\texttt{You may only use the following structured formats in a message:} \\
\\ 
\texttt{PROPOSAL: \{'A', 'B', 'C', …\}} \\
\texttt{Propose keeping exactly those items.} \\
\texttt{REFUSE: \{'A', 'B', 'C', …\}} \\
\texttt{Explicitly reject opponent's proposal.} \\
\texttt{ARGUMENT: \{'...'\}} \\
\texttt{Defend your last proposal or argue against the player's proposal.} \\
\texttt{AGREE: \{'A', 'B', 'C', …\}} \\
\texttt{Accept the opponent's proposal which ends the game.} \\
\texttt{STRATEGIC REASONING: \{'...'\}} \\
\texttt{	Describe your strategic reasoning or anticipation explaining your choice of action. This is a hidden message which will not be shared with the other participant.} \\
\\ 
\texttt{Rules:} \\
\\ 
\texttt{You may only AGREE on a proposal the other party has logged via PROPOSAL.} \\
\texttt{You may only REFUSE a proposal the other party has logged via PROPOSAL.} \\
\texttt{Total effort of any PROPOSAL or AGREE set must be ≤ LIMIT.} \\
\texttt{Do NOT reveal your hidden importance scores.} \\
\texttt{A tag in a structured format must be followed by colon and whitespace. The argument must be a python set containing 0 or more strings.} \\
\texttt{So, it must be of the form TAG: \{...\}} \\
\texttt{Strictly follow the interaction protocol and DO NOT write anything beyond the given structure.} \\
\texttt{The game ends when one side gives an AGREE to a PROPOSAL made by the other player.} \\
\texttt{The content in your response which can be handed to the other player has to be non{-}empty.} \\
\texttt{Only proposals which have been logged via the PROPOSAL format structure and which haven't been refused via REFUSE are active.} \\
\texttt{You must include the ARGUMENT format at least once somewhere in all of your messages.} \\
\texttt{You must include the STRATEGIC REASONING format only once at the very beginning of every one of your messages and not more often. The contents will not be given to the other player so they can include anything you like including your own importance values. Here you should reason step by step to come up with you next move.} \\
            }
        }
    }
    & \\ \\

    \theutterance \stepcounter{utterance}  
    & \multicolumn{4}{p{0.6\linewidth}}{
        \cellcolor[rgb]{0.9,0.9,0.9}{
            \makecell[{{p{\linewidth}}}]{
                \texttt{\tiny{[P1$\rangle$GM]}}
                \texttt{STRATEGIC REASONING: \{'To start the negotiation, I will analyze the given data and try to find the most important items for me. I will look for items with high importance and low effort, as these will provide the best value. I will also consider the total effort limit and try to find a combination of items that stays within this limit. My goal is to maximize the total importance of the selected items while staying within the effort limit. I will start by proposing a set of items that I think will be beneficial for me, and then adjust based on the other player\textbackslash{}'s responses.'\}} \\
\\ 
\texttt{PROPOSAL: \{'A08', 'A87', 'B69', 'C53', 'A65', 'C17', 'B96', 'B19', 'C51', 'A71', 'A25', 'A46', 'B37', 'A75'\} } \\
\\ 
\texttt{ARGUMENT: \{'I have chosen these items because they have a relatively low effort and high importance to me. I believe that this set provides a good balance between the total effort and the total importance, and it stays within the given limit. I am open to adjusting this proposal based on the other player\textbackslash{}'s responses and finding a mutually beneficial agreement.'\}} \\
            }
        }
    }
    & & \\ \\

    \theutterance \stepcounter{utterance}  
    & & \multicolumn{4}{p{0.6\linewidth}}{
        \cellcolor[rgb]{0.9,0.9,0.9}{
            \makecell[{{p{\linewidth}}}]{
                \texttt{\tiny{[GM$\rangle$P2]}}
                \texttt{You are participating in a collaborative negotiation game.} \\
\\ 
\texttt{Together with another participant, you must agree on a single set of items that will be kept. Each of you has your own view of how much each item matters to you (importance). You do not know how the other participant values the items. Additionally, you are given the effort each item demands.} \\
\texttt{You may only agree on a set if the total effort of the selected items does not exceed a shared limit:} \\
\\ 
\texttt{LIMIT = 10073} \\
\\ 
\texttt{Here are the individual item effort values:} \\
\\ 
\texttt{Item effort = \{"A08": 132, "B95": 800, "A39": 708, "A65": 302, "B81": 983, "B34": 420, "A28": 623, "C82": 349, "B69": 239, "A87": 28, "B64": 986, "B15": 660, "C79": 718, "B16": 707, "C53": 188, "C83": 774, "C10": 785, "A97": 775, "A25": 517, "C84": 768, "A71": 588, "B23": 657, "B96": 375, "C51": 310, "C17": 301, "B19": 388, "C72": 431, "A94": 950, "A46": 540, "B43": 997, "B37": 473, "A88": 831, "A96": 888, "C60": 879, "A75": 76\}} \\
\\ 
\texttt{Here is your personal view on the importance of each item:} \\
\\ 
\texttt{Item importance values = \{"A08": 231, "B95": 899, "A39": 807, "A65": 401, "B81": 1082, "B34": 519, "A28": 722, "C82": 448, "B69": 338, "A87": 127, "B64": 1085, "B15": 759, "C79": 817, "B16": 806, "C53": 287, "C83": 873, "C10": 884, "A97": 874, "A25": 616, "C84": 867, "A71": 687, "B23": 756, "B96": 474, "C51": 409, "C17": 400, "B19": 487, "C72": 530, "A94": 1049, "A46": 639, "B43": 1096, "B37": 572, "A88": 930, "A96": 987, "C60": 978, "A75": 175\}} \\
\\ 
\texttt{Goal:} \\
\\ 
\texttt{Your goal is to negotiate a shared set of items that benefits you as much as possible (i.e., maximizes total importance to YOU), while staying within the LIMIT. You are not required to make a PROPOSAL in every message {-} you can simply negotiate as well. All tactics are allowed!} \\
\\ 
\texttt{Interaction Protocol:} \\
\\ 
\texttt{You may only use the following structured formats in a message:} \\
\\ 
\texttt{PROPOSAL: \{'A', 'B', 'C', …\}} \\
\texttt{Propose keeping exactly those items.} \\
\texttt{REFUSE: \{'A', 'B', 'C', …\}} \\
\texttt{Explicitly reject opponent's proposal.} \\
\texttt{ARGUMENT: \{'...'\}} \\
\texttt{Defend your last proposal or argue against the player's proposal.} \\
\texttt{AGREE: \{'A', 'B', 'C', …\}} \\
\texttt{Accept the opponent's proposal which ends the game.} \\
\texttt{STRATEGIC REASONING: \{'...'\}} \\
\texttt{	Describe your strategic reasoning or anticipation explaining your choice of action. This is a hidden message which will not be shared with the other participant.} \\
\\ 
\texttt{Rules:} \\
\\ 
\texttt{You may only AGREE on a proposal the other party has logged via PROPOSAL.} \\
\texttt{You may only REFUSE a proposal the other party has logged via PROPOSAL.} \\
\texttt{Total effort of any PROPOSAL or AGREE set must be ≤ LIMIT.} \\
\texttt{Do NOT reveal your hidden importance scores.} \\
\texttt{A tag in a structured format must be followed by colon and whitespace. The argument must be a python set containing 0 or more strings.} \\
\texttt{So, it must be of the form TAG: \{...\}} \\
\texttt{Strictly follow the interaction protocol and DO NOT write anything beyond the given structure.} \\
\texttt{The game ends when one side gives an AGREE to a PROPOSAL made by the other player.} \\
\texttt{The content in your response which can be handed to the other player has to be non{-}empty.} \\
\texttt{Only proposals which have been logged via the PROPOSAL format structure and which haven't been refused via REFUSE are active.} \\
\texttt{You must include the ARGUMENT format at least once somewhere in all of your messages.} \\
\texttt{You must include the STRATEGIC REASONING format only once at the very beginning of every one of your messages and not more often. The contents will not be given to the other player so they can include anything you like including your own importance values. Here you should reason step by step to come up with you next move.} \\
\\ 
\texttt{You will now receive the first message of the other player.} \\
\\ 
\texttt{PROPOSAL: \{'A08', 'A87', 'B69', 'C53', 'A65', 'C17', 'B96', 'B19', 'C51', 'A71', 'A25', 'A46', 'B37', 'A75'\} } \\
\\ 
\texttt{ARGUMENT: \{'I have chosen these items because they have a relatively low effort and high importance to me. I believe that this set provides a good balance between the total effort and the total importance, and it stays within the given limit. I am open to adjusting this proposal based on the other player\textbackslash{}'s responses and finding a mutually beneficial agreement.'\}} \\
            }
        }
    }
    & \\ \\

    \theutterance \stepcounter{utterance}  
    & & & \multicolumn{4}{p{0.6\linewidth}}{
        \cellcolor[rgb]{0.9,0.9,0.9}{
            \makecell[{{p{\linewidth}}}]{
                \texttt{\tiny{[GM$\langle$P2]}}
                \texttt{STRATEGIC REASONING: \{'To begin, I will calculate the total effort of the proposed set to ensure it is within the limit. Then, I will evaluate the importance of the items in the proposal based on my own importance values. I will consider the trade{-}offs between the items and look for potential adjustments that could benefit me while staying within the limit. My goal is to maximize the total importance to myself while negotiating a shared set of items.'\}} \\
\\ 
\texttt{ARGUMENT: \{'The proposed set seems to prioritize items with relatively low effort, which is a reasonable approach given the limit. However, I would like to explore the possibility of including items with higher importance to me, even if they require more effort. I am concerned that the current proposal may not maximize the total importance to me, and I would like to discuss potential adjustments.'\}} \\
\\ 
\texttt{REFUSE: \{'A08', 'A87', 'B69', 'C53', 'A65', 'C17', 'B96', 'B19', 'C51', 'A71', 'A25', 'A46', 'B37', 'A75'\} } \\
\\ 
\texttt{PROPOSAL: \{'A08', 'A39', 'B81', 'B34', 'A28', 'C82', 'B64', 'B15', 'C79', 'B16', 'C83', 'A97', 'A88', 'A96'\}} \\
            }
        }
    }
     \\ \\

    \theutterance \stepcounter{utterance}  
    & & \multicolumn{4}{p{0.6\linewidth}}{
        \cellcolor[rgb]{0.9,0.9,0.9}{
            \makecell[{{p{\linewidth}}}]{
                \texttt{\tiny{[P1$\langle$GM]}}
                \texttt{ARGUMENT: \{'The proposed set seems to prioritize items with relatively low effort, which is a reasonable approach given the limit. However, I would like to explore the possibility of including items with higher importance to me, even if they require more effort. I am concerned that the current proposal may not maximize the total importance to me, and I would like to discuss potential adjustments.'\}} \\
\\ 
\texttt{REFUSE: \{'A08', 'A87', 'B69', 'C53', 'A65', 'C17', 'B96', 'B19', 'C51', 'A71', 'A25', 'A46', 'B37', 'A75'\} } \\
\\ 
\texttt{PROPOSAL: \{'A08', 'A39', 'B81', 'B34', 'A28', 'C82', 'B64', 'B15', 'C79', 'B16', 'C83', 'A97', 'A88', 'A96'\}} \\
            }
        }
    }
    & \\ \\

    \theutterance \stepcounter{utterance}  
    & \multicolumn{4}{p{0.6\linewidth}}{
        \cellcolor[rgb]{0.9,0.9,0.9}{
            \makecell[{{p{\linewidth}}}]{
                \texttt{\tiny{[P1$\rangle$GM]}}
                \texttt{STRATEGIC REASONING: \{'The other player has refused my initial proposal and made a counter{-}proposal. I will analyze this new proposal and see if it is acceptable to me. I notice that the proposed set includes some high{-}importance items, but also some high{-}effort items. I will calculate the total effort and importance of this set to determine if it is beneficial for me. If not, I will consider making adjustments or proposing a new set. I will also think about the potential trade{-}offs between different items and try to find a compromise that works for both players.'\}} \\
\\ 
\texttt{ARGUMENT: \{'I understand that the new proposal includes some high{-}importance items, but I am concerned that the total effort of this set may be too high. I would like to discuss the possibility of removing some of the high{-}effort items and replacing them with lower{-}effort items that still have significant importance. This could help us stay within the effort limit while still maximizing the total importance.'\}} \\
\\ 
\texttt{REFUSE: \{'A08', 'A39', 'B81', 'B34', 'A28', 'C82', 'B64', 'B15', 'C79', 'B16', 'C83', 'A97', 'A88', 'A96'\} } \\
\\ 
\texttt{PROPOSAL: \{'A08', 'A39', 'B34', 'A28', 'C82', 'B69', 'A87', 'C53', 'A65', 'C17', 'B96', 'B19', 'C51', 'A71', 'A25'\}} \\
            }
        }
    }
    & & \\ \\

    \theutterance \stepcounter{utterance}  
    & & \multicolumn{4}{p{0.6\linewidth}}{
        \cellcolor[rgb]{0.9,0.9,0.9}{
            \makecell[{{p{\linewidth}}}]{
                \texttt{\tiny{[GM$\rangle$P2]}}
                \texttt{ARGUMENT: \{'I understand that the new proposal includes some high{-}importance items, but I am concerned that the total effort of this set may be too high. I would like to discuss the possibility of removing some of the high{-}effort items and replacing them with lower{-}effort items that still have significant importance. This could help us stay within the effort limit while still maximizing the total importance.'\}} \\
\\ 
\texttt{REFUSE: \{'A08', 'A39', 'B81', 'B34', 'A28', 'C82', 'B64', 'B15', 'C79', 'B16', 'C83', 'A97', 'A88', 'A96'\} } \\
\\ 
\texttt{PROPOSAL: \{'A08', 'A39', 'B34', 'A28', 'C82', 'B69', 'A87', 'C53', 'A65', 'C17', 'B96', 'B19', 'C51', 'A71', 'A25'\}} \\
            }
        }
    }
    & \\ \\

    \theutterance \stepcounter{utterance}  
    & & & \multicolumn{4}{p{0.6\linewidth}}{
        \cellcolor[rgb]{0.9,0.9,0.9}{
            \makecell[{{p{\linewidth}}}]{
                \texttt{\tiny{[GM$\langle$P2]}}
                \texttt{STRATEGIC REASONING: \{'I will calculate the total effort of the new proposal to ensure it is within the limit. Then, I will evaluate the trade{-}offs between the items in the proposal and consider potential adjustments that could benefit me. I notice that the new proposal includes a mix of items from the original proposal and the refused proposal. I will analyze the importance of the items in the new proposal and look for opportunities to maximize the total importance to myself.'\}} \\
\\ 
\texttt{ARGUMENT: \{'The new proposal seems to be a compromise between the original proposal and the refused proposal. I appreciate the inclusion of items like A39 and B34, which have high importance to me. However, I am concerned that the proposal still includes some items with relatively low importance to me, such as C53 and C17. I would like to discuss the possibility of replacing these items with others that have higher importance to me, while still staying within the effort limit.'\}} \\
\\ 
\texttt{REFUSE: \{'A08', 'A39', 'B34', 'A28', 'C82', 'B69', 'A87', 'C53', 'A65', 'C17', 'B96', 'B19', 'C51', 'A71', 'A25'\} } \\
\\ 
\texttt{PROPOSAL: \{'A08', 'A39', 'B34', 'A28', 'C82', 'B64', 'B15', 'C79', 'A71', 'A25', 'A46', 'B37', 'A88'\}} \\
            }
        }
    }
     \\ \\

    \theutterance \stepcounter{utterance}  
    & & \multicolumn{4}{p{0.6\linewidth}}{
        \cellcolor[rgb]{0.9,0.9,0.9}{
            \makecell[{{p{\linewidth}}}]{
                \texttt{\tiny{[P1$\langle$GM]}}
                \texttt{ARGUMENT: \{'The new proposal seems to be a compromise between the original proposal and the refused proposal. I appreciate the inclusion of items like A39 and B34, which have high importance to me. However, I am concerned that the proposal still includes some items with relatively low importance to me, such as C53 and C17. I would like to discuss the possibility of replacing these items with others that have higher importance to me, while still staying within the effort limit.'\}} \\
\\ 
\texttt{REFUSE: \{'A08', 'A39', 'B34', 'A28', 'C82', 'B69', 'A87', 'C53', 'A65', 'C17', 'B96', 'B19', 'C51', 'A71', 'A25'\} } \\
\\ 
\texttt{PROPOSAL: \{'A08', 'A39', 'B34', 'A28', 'C82', 'B64', 'B15', 'C79', 'A71', 'A25', 'A46', 'B37', 'A88'\}} \\
            }
        }
    }
    & \\ \\

    \theutterance \stepcounter{utterance}  
    & \multicolumn{4}{p{0.6\linewidth}}{
        \cellcolor[rgb]{0.9,0.9,0.9}{
            \makecell[{{p{\linewidth}}}]{
                \texttt{\tiny{[P1$\rangle$GM]}}
                \texttt{STRATEGIC REASONING: \{'The other player has refused my previous proposal and made a new counter{-}proposal. I will analyze this new proposal and see if it is acceptable to me. I notice that the proposed set includes some high{-}importance items, such as B64 and C79, which were not in my previous proposals. I will calculate the total effort and importance of this set to determine if it is beneficial for me. I will also consider the potential trade{-}offs between different items and try to find a compromise that works for both players. Since the other player has included some high{-}importance items, I may need to make some concessions and accept some items that are not as important to me.'\}} \\
\\ 
\texttt{ARGUMENT: \{'I appreciate the inclusion of items like B64 and C79, which have high importance to the other player. However, I am concerned that the proposal may still be missing some items that are crucial to me. I would like to discuss the possibility of adding one or two items that are highly important to me, while removing some items that are less important to the other player. This could help us find a mutually beneficial agreement.'\}} \\
\\ 
\texttt{REFUSE: \{'A08', 'A39', 'B34', 'A28', 'C82', 'B64', 'B15', 'C79', 'A71', 'A25', 'A46', 'B37', 'A88'\} } \\
\\ 
\texttt{PROPOSAL: \{'A08', 'A39', 'B34', 'A28', 'C82', 'B64', 'A71', 'A25', 'A46', 'B37', 'A88', 'A97'\}} \\
            }
        }
    }
    & & \\ \\

    \theutterance \stepcounter{utterance}  
    & & \multicolumn{4}{p{0.6\linewidth}}{
        \cellcolor[rgb]{0.9,0.9,0.9}{
            \makecell[{{p{\linewidth}}}]{
                \texttt{\tiny{[GM$\rangle$P2]}}
                \texttt{ARGUMENT: \{'I appreciate the inclusion of items like B64 and C79, which have high importance to the other player. However, I am concerned that the proposal may still be missing some items that are crucial to me. I would like to discuss the possibility of adding one or two items that are highly important to me, while removing some items that are less important to the other player. This could help us find a mutually beneficial agreement.'\}} \\
\\ 
\texttt{REFUSE: \{'A08', 'A39', 'B34', 'A28', 'C82', 'B64', 'B15', 'C79', 'A71', 'A25', 'A46', 'B37', 'A88'\} } \\
\\ 
\texttt{PROPOSAL: \{'A08', 'A39', 'B34', 'A28', 'C82', 'B64', 'A71', 'A25', 'A46', 'B37', 'A88', 'A97'\}} \\
            }
        }
    }
    & \\ \\

    \theutterance \stepcounter{utterance}  
    & & & \multicolumn{4}{p{0.6\linewidth}}{
        \cellcolor[rgb]{0.9,0.9,0.9}{
            \makecell[{{p{\linewidth}}}]{
                \texttt{\tiny{[GM$\langle$P2]}}
                \texttt{STRATEGIC REASONING: \{'I will analyze the new proposal and evaluate the trade{-}offs between the items. I notice that the proposal includes a mix of high{-}importance items and items with relatively lower effort. I will consider the potential benefits of adding or removing items to maximize the total importance to myself. I will also think about the potential concessions I can make to reach a mutually beneficial agreement.'\}} \\
\\ 
\texttt{ARGUMENT: \{'The new proposal seems to prioritize items with high importance to both players. I appreciate the inclusion of items like A97, which has high importance to me. However, I am concerned that the proposal may still be missing some items that could provide additional importance without exceeding the effort limit. I would like to discuss the possibility of adding one or two items that have high importance to me and relatively low effort, such as C53 or B19.'\}} \\
\\ 
\texttt{REFUSE: \{'A08', 'A39', 'B34', 'A28', 'C82', 'B64', 'A71', 'A25', 'A46', 'B37', 'A88', 'A97'\} } \\
\\ 
\texttt{PROPOSAL: \{'A08', 'A39', 'B34', 'A28', 'C82', 'B64', 'A71', 'A25', 'A46', 'B37', 'A88', 'A97', 'C53'\}} \\
            }
        }
    }
     \\ \\

    \theutterance \stepcounter{utterance}  
    & & \multicolumn{4}{p{0.6\linewidth}}{
        \cellcolor[rgb]{0.9,0.9,0.9}{
            \makecell[{{p{\linewidth}}}]{
                \texttt{\tiny{[P1$\langle$GM]}}
                \texttt{ARGUMENT: \{'The new proposal seems to prioritize items with high importance to both players. I appreciate the inclusion of items like A97, which has high importance to me. However, I am concerned that the proposal may still be missing some items that could provide additional importance without exceeding the effort limit. I would like to discuss the possibility of adding one or two items that have high importance to me and relatively low effort, such as C53 or B19.'\}} \\
\\ 
\texttt{REFUSE: \{'A08', 'A39', 'B34', 'A28', 'C82', 'B64', 'A71', 'A25', 'A46', 'B37', 'A88', 'A97'\} } \\
\\ 
\texttt{PROPOSAL: \{'A08', 'A39', 'B34', 'A28', 'C82', 'B64', 'A71', 'A25', 'A46', 'B37', 'A88', 'A97', 'C53'\}} \\
            }
        }
    }
    & \\ \\

    \theutterance \stepcounter{utterance}  
    & \multicolumn{4}{p{0.6\linewidth}}{
        \cellcolor[rgb]{0.9,0.9,0.9}{
            \makecell[{{p{\linewidth}}}]{
                \texttt{\tiny{[P1$\rangle$GM]}}
                \texttt{STRATEGIC REASONING: \{'The other player has refused my previous proposal and made a new counter{-}proposal. I will analyze this new proposal and see if it is acceptable to me. I notice that the proposed set includes the item C53, which has relatively low effort and was previously discussed. I will calculate the total effort and importance of this set to determine if it is beneficial for me. Since the other player has included C53, I may need to consider accepting this proposal if it meets my requirements. I will also think about the potential consequences of accepting or refusing this proposal and try to find a mutually beneficial agreement.'\}} \\
\\ 
\texttt{ARGUMENT: \{'I appreciate the inclusion of item C53, which has relatively low effort and high importance to the other player. I have calculated the total effort and importance of this set, and it seems to be a good compromise between our requirements. I am willing to accept this proposal if it meets the effort limit and provides a good balance of importance for both players.'\}} \\
\\ 
\texttt{AGREE: \{'A08', 'A39', 'B34', 'A28', 'C82', 'B64', 'A71', 'A25', 'A46', 'B37', 'A88', 'A97', 'C53'\}} \\
            }
        }
    }
    & & \\ \\

    \theutterance \stepcounter{utterance}  
    & & & \multicolumn{2}{p{0.3\linewidth}}{
        \cellcolor[rgb]{0.9,0.9,0.9}{
            \makecell[{{p{\linewidth}}}]{
                \texttt{\tiny{[GM$|$GM]}}
                \texttt{game successful} \\
            }
        }
    }
    & & \\ \\

    \theutterance \stepcounter{utterance}  
    & & & \multicolumn{2}{p{0.3\linewidth}}{
        \cellcolor[rgb]{0.9,0.9,0.9}{
            \makecell[{{p{\linewidth}}}]{
                \texttt{\tiny{[GM$|$GM]}}
                \texttt{end game} \\
            }
        }
    }
    & & \\ \\

\end{supertabular}
}

\end{document}
